\begin{frame}{연습문제(15.2)}
  \small
  \begin{enumerate}
    \item \textbf{(수치 계산)} 부모 방문 $N(s)=111$, 자식
      $A: (Q,N)=(10,100)$, $B: (5,10)$, $C: (0,1)$, $c=1.4$일
      때 세 자식의 UCT 값 계산 후 어떤 수가 선택되는지
      판정. (힌트: $\ln 111 \approx 4.71$. ``승률 0인 C''가
      \textbf{다시} 선택될 수 있을까?)
    \item \textbf{(개념 서술)} 선택 단계의 정지 조건
      \texttt{len(node.children) == len(nim\_moves(...))}가
      없으면(\texttt{while node.children:}으로 대체) 어떤
      증상이 나서 \textbf{왜} 그런지 ``확장의 보장''과
      연결지어. (힌트: 버그 투어 1의 재현 값 인용.)
    \item \textbf{(코딩)} \texttt{mcts\_nim}에서 $c$를 0으로
      바꿔 3000회 실행 --- 방문 분포가 어떻게 되고,
      \textbf{어느 자식}에 집중되는지 예측(힌트: 어떤 수를
      \textbf{먼저} 확장하는지는 \texttt{rng} 순서에 의존) 후
      확인. ``탐험 항 $c\sqrt{\ln N / N}$의 역할''에 대해
      무엇을 말하는가?
    \item \textbf{(개념 서술)} 님(5)에서 이기는 수의 MCTS
      추정 승률(0.987)이 ``무작위 상대 참 승률''(2/3)보다 큰
      이유를, ``트리가 기억하는 부분''과 ``롤아웃이
      표본추출하는 부분''의 \textbf{분업}으로. 이 갭은 어떤
      게임에서 커지고, 어떤 게임에서 작으며, 왜?
      (힌트: 남은 게임의 \textbf{깊이}.)
  \end{enumerate}
\end{frame}
